A continuaci\'on se describe el flujo paso a paso tal como ocurre en la actualidad:

\begin{enumerate}
    \item \textbf{Configuraci\'on previa (Comisi\'on de Matr\'icula):} Antes de abrir el
    per\'iodo de matr\'icula, la comisi\'on registra en el SIGA-ORCE la oferta de cursos, los
    turnos asignados por facultad o especialidad, y los cupos m\'aximos por secci\'on.

    \item \textbf{Acceso al sistema (Estudiante):} El estudiante ingresa a la Intranet
    DIRCE con su c\'odigo de matr\'icula y contrase\~na institucional. El sistema autentica sus
    credenciales y verifica si su turno de matr\'icula se encuentra vigente.

    \item \textbf{Despliegue de oferta (Sistema):} Una vez validado el turno, el sistema
    presenta la lista de cursos habilitados para el semestre, con informaci\'on de horarios,
    secciones y vacantes disponibles en tiempo real.

    \item \textbf{Selecci\'on de cursos (Estudiante):} El estudiante elige los cursos de
    su plan de estudios y construye un horario tentativo. Esta actividad puede repetirse
    m\'ultiples veces si las validaciones posteriores detectan inconsistencias.

    \item \textbf{Validaci\'on de cr\'editos (Sistema):} Conforme al Reglamento RR-0570-2022, el sistema calcula la suma de cr\'editos de los cursos seleccionados y la compara con el l\'imite permitido seg\'un el promedio ponderado acumulado del estudiante.
    Si el l\'imite es superado, el estudiante debe ajustar su selecci\'on.

    \item \textbf{Validaci\'on de horarios (Sistema):} Se verifica que no existan cruces
    entre los horarios de los cursos seleccionados. Si se detecta un conflicto, el
    estudiante regresa a la etapa de selecci\'on.

    \item \textbf{Validaci\'on de vacantes (Sistema):} El sistema confirma que al menos
    una vacante est\'e disponible en cada secci\'on elegida. Si no hay vacantes, el caso es
    derivado autom\'aticamente al Subproceso de Rectificaci\'on (SP3A), donde el estudiante
    podr\'a solicitar un cupo adicional.

    \item \textbf{Confirmaci\'on de inscripci\'on (Estudiante):} Superadas las tres
    validaciones, el estudiante revisa su horario final y confirma la inscripci\'on.

    \item \textbf{Emisi\'on de Pre-Ficha (Sistema):} El sistema registra formalmente la
    matr\'icula, genera la Pre-Ficha con el detalle de cursos y cr\'editos inscritos, y
    reserva las vacantes seleccionadas para el estudiante.
\end{enumerate}

\noindent
\textbf{Problem\'atica identificada en el proceso AS IS:} El flujo actual presenta
dependencia total de la conectividad a la Intranet DIRCE. Ante ca\'idas del sistema
durante los turnos de matr\'icula, el proceso queda bloqueado sin procedimiento
alternativo documentado. La ausencia de notificaciones proactivas obliga al
estudiante a monitorear manualmente la disponibilidad de vacantes, generando reingresos
repetidos al sistema. Estos problemas se abordan en la propuesta de mejora (proceso
TO BE) del Subproceso~2.
